\documentclass[11pt,a4paper]{article}
\usepackage[ngerman]{babel}
\usepackage[utf8]{inputenc}
\usepackage[T1]{fontenc}
\usepackage{lmodern}
\usepackage{csquotes}
\usepackage{amsmath,amssymb,amsthm}
\usepackage{booktabs}
\usepackage{geometry}
\geometry{margin=2.4cm}
\usepackage[hidelinks]{hyperref}
\usepackage{enumitem}
\usepackage{microtype}
\newtheorem{satz}{Satz}
\newtheorem{lemma}[satz]{Lemma}
\newcommand{\Om}{\Omega}
\newcommand{\Tr}{\operatorname{Tr}}
\newcommand{\tr}{\operatorname{tr}}
\title{TFPT-Anschluss: Zustandsgate, Trägerzellen und der $A_3$-Faktor des Seams\\
\large Ergebnisse der Fable-Runden 3 (11.--14.\ September 2026) als Fortsetzung des Manuskripts
\enquote{Vom geometrischen Compiler zum Prozess mit Aufzeichnung} (Version 1.1)}
\author{Claude Fable 5.1 im Cursor-Fenster \texttt{tfpt-theoryv4}\\
\small für Stefan Hamann\\ \small Prüfer und Zahlen: \texttt{experiments/theory-contracts/}\\ \small\texttt{universalraum-fable-kernfragen-20260910/fable-runde3/}}
\date{14.\ September 2026 · Forschungsmanuskript, kein begutachteter Abschluss}

